\styledchapter{Independence}
\section{Preliminaries}

Consider a measure space $\left( \Omega, \mathscr{F}, \P\right)$. 

\begin{notation}
	For a random variable $X: \left( \Omega, \mathscr{F}, \P \right) \to \left( \R, \mathscr{B}_{\R} \right)$, when we write $\{X \in B\}$, we are really referring to the set as \[
		\{X \in B\} = \{\omega \in \Omega: X(\omega) \in B\}
	.\] Similarly, \[
		\{X \le a\} = \{\omega \in \Omega: X(\omega) \le a\}
	,\] \[
		\{X \in A, Y \in B\} = \{\omega \in \Omega: X(\omega) \in A, Y(\omega) \in B\}
	,\] and so on.

	Perhaps most importantly, \[
		\P\left[X \in B\right] = \P\left[\omega \in \Omega: X(\omega) \in B\right] 
	.\] 
\end{notation}

\begin{definition}[Independent events]
	We say that the events $A_1,\ldots,A_n \in \mathscr{F}$ are independent if for every $\Lambda \subseteq \{1,\ldots,n\}$, \[
		\P \left[ \bigcap\limits_{i\in \Lambda}^{} A_i \right] = \prod_{i \in \Lambda}^{} \P\left[A_i\right]  
	.\] 

	We say $\mathscr{E}_1,\ldots,\mathscr{E}_n \subseteq \mathscr{F}$ are independent if for any $A_1 \in \mathscr{E}_1,\ldots,A_n \in \mathscr{E}_n$, $\left( A_1,\ldots,A_n \right)$ are independent.
\end{definition}

\begin{definition}[Generated $\sigma$-field (of a random variable)]
	Let $X: (\Omega, \mathscr{F}, \P) \to \left( \R,\mathscr{B}_{\R} \right)$ be a random variable. We define $\sigma(X)$ to be the smallest $\sigma$-field such that \[
		X: (\Omega, \sigma(X), \P) \to \left( \R,\mathscr{B}_\R \right)
	\] is measurable.

	Since $X ^{-1} \mathscr{B}_{\R}$ is a $\sigma$-field containing $X$ (preimage of a $\sigma$-field is also a $\sigma$-field), \[
		\sigma(X) \subseteq  X^{-1}\mathscr{B}_\R = X^{-1} \sigma\left( \{\{x > a\}: a \in \Q\} \right)
	.\]
	Conversely, suppose
	\[
		\begin{aligned}
			X: \left( \Omega, \sigma\left( X \right), \P \right)
			&\to \left( \R, \mathscr{B}_\R \right)
		\end{aligned}
	\]
	is measurable.
	Then $X ^{-1}(B) \in \sigma\left( X \right)$ for every $B \in \mathscr{B}_\R$.
	Hence,
	\[
		X ^{-1}\left( \mathscr{B}_\R \right) \subseteq \sigma\left( X \right)
	.\]
	Therefore,
	\[
		\sigma(X) = X ^{-1}\mathscr{B}_\R
	.\] 
\end{definition}

\begin{example}
	Suppose $A \in \mathscr{F}$ and $X = \mathds{1}_{A}$. Then \[
		\sigma(X) = \{\O, \Omega, A, A^{c}\}
	.\] 
\end{example}

\begin{definition}[Independent random variables]
	We say that $X_1,\ldots,X_n: (\Omega, \mathscr{F}, \P) \to \left( \R,\mathscr{B}_\R \right)$ are independent if $\sigma\left( X_1 \right),\ldots,\sigma(X_n)$ are independent.
\end{definition}

Why does this definition make sense? 
For $X \ind Y$ and for any $A,B \in \mathscr{B}_\R$, we want \[
	\P\left[X \in A, Y \in B\right] = \P\left[X \in A\right]  \P\left[Y \in B\right] 
.\] We see that \[
	\begin{aligned}
		\P\left[X \in A\right]  &= \P\left[X ^{-1}A\right], \\
		\P\left[Y \in B\right] &= \P\left[Y ^{-1} B\right], \\
		\P\left[X \in A, Y \in B\right]  &= \P\left[X ^{-1}A \cap Y ^{-1}B\right]
	\end{aligned}
.\] We need $\P$ to be defined on and satisfy the independence condition on sets of this last type.

\begin{lemma}
	If $X \ind Y$ and $\E \left| X \right|,\E \left| Y \right| < \infty$, then \[
		\E\left[XY\right]  = \E\left[X\right] \E\left[Y\right] 
	.\] 
\end{lemma}
\begin{proof}
	\begin{enumerate}[label=(\roman*)]
		\item If $X = \mathds{1}_{A}, Y = \mathds{1}_{B}$, then $XY = \mathds{1}_{A\cap B}$. Hence
			\begin{align*}
				\E\left[XY\right] &= \E\left[\mathds{1}_{A\cap B}\right] \\
				&= \P\left[A \cap B\right]  \\
				&=\P\left[A\right] \P\left[B\right]\text{ by independence} \\
				&= \E\left[X\right] \E\left[Y\right] 
			.\end{align*}
		\item For $\{A_i\}, \{B_j\}$ partitions of $\Omega$, if $X = \sum_i a_i \mathds{1}_{A_i}$ and $Y = \sum_j b_j \mathds{1}_{B_j}$, then we are tempted to write
			\begin{align*}
				XY &=  \sum\limits_{i=1}^{n} \sum\limits_{j=1}^{m} a_i b_j \mathds{1}_{A_i \cap B_j} \\
				\E\left[XY\right] &= \sum\limits_{i=1}^{n} \sum\limits_{j=1}^{m} a_i b_j \P\left[A _i \cap B_j\right]  \\
				&= \sum\limits_{i=1}^{n} \sum\limits_{j=1}^{m} a_i b_j \P\left[A_i\right] \P\left[B_j\right]
			\end{align*}
			and continue, but we don't know $A_i$ and $B_j$ are independent, since this requires $A_i \in \sigma(X)$ and $B_j \in \sigma(Y)$. This is not guaranteed if the $a_i$, $b_j$ are not unique. (For example, if $X(\omega) = 1$ for $\omega \in A_1$ and $\omega \in A_2$, then $A_1 \not\in \sigma(X)$, since $A_1$ being in $\sigma(X)$ is not necessary for $X$ to be measurable --- only $A_1 \cup A_2 \in \sigma\left( X \right)$ is required.
			
			Instead, we should decompose as \[
				X = \sum\limits_{i=1}^{n} a_i \mathds{1}_{X = a_i}, \quad Y = \sum\limits_{j=1}^{m} b_j \mathds{1}_{Y = b_j}
			.\] Then we can actually continue to get
			\begin{align*}
				\E\left[XY\right]  &=  \sum\limits_{i=1}^{n} \sum\limits_{j=1}^{m} a_i b_j \P\left[X = a_i, Y = b_j\right]  \\
				&= \sum\limits_{i=1}^{n} \sum\limits_{j=1}^{m} a_i b_j \P\left[X = a_i\right] \P\left[Y = b_j\right]   \\
				&= \left( \sum\limits_{i=1}^{n} a_i \P\left[X = a_i\right]  \right)\left( \sum\limits_{j=1}^{m} b_j \P\left[Y = b_j\right]  \right) \\
				&= \E\left[X\right] \E\left[Y\right] 
			.\end{align*}
		\item Now consider $X,Y \ge 0$. Write 
			\begin{align*}
				X_n &= \sum\limits_{k=0}^{n2^n - 1} \frac{k}{2^n}\mathds{1}_{\frac{k}{2^n}\le X < \frac{k+1}{2^n}} + n \mathds{1}_{X \ge n} \\
				Y_n &= \sum\limits_{k=0}^{n2^n - 1} \frac{k}{2^n}\mathds{1}_{\frac{k}{2^n}\le Y < \frac{k+1}{2^n}} + n \mathds{1}_{Y \ge n}
			.\end{align*}
			Then $X_n \uparrow X$, $Y_n \uparrow Y$.
			Since $\sigma(X_n) \subseteq \sigma(X)$\boxednote{Consider $X = 3\mathds{1}_{A_1} + 4\mathds{1}_{A_2} - \mathds{1}_{A_3}$ for disjoint $A_1,A_2,A_3$, and find $\sigma(X)$. This provides intuition for what $\sigma(X)$ generally is.} and $\sigma(Y_n) \subseteq \sigma(Y)$ imply that $X_n \ind Y_n$, \[
				\E\left[X_n Y_n\right]  = \E\left[X_n\right] \E\left[Y_n\right] 
			.\] Taking $n\to \infty$ concludes.
		\item For general $X,Y$, put $X = X^+ - X^-$ and $Y = Y^+ - Y^-$, so \[
			XY = X^+Y^+ + X^-Y^- - X^+Y^- - X^-Y^+
		.\] Since $\sigma(X^+) \subseteq \sigma(X)$ and so on, we see that each term in the sum satisfies independence; we can take the expectation termwise.
		Using linearity of the integral to recombine $X^+$ with $X^-$ and similarly for $Y$, we get independence of $X$ and $Y$.
	\end{enumerate}
\end{proof}

\begin{definition}[Identically distributed random variables]
	We say that $X_1,\ldots,X_n$ are identically distributed if \[
		\P\left[X_j \in A\right] = \P\left[X_i \in A\right] 
	\] for all $i,j$ and $A \in \mathscr{F}$.
\end{definition}
